\section{Consequences and prediction conventions}

\begin{corollary}[Individual posteriors and clipped draws]
For the law in Theorem~\ref{thm:relu-mixture-failure}, at all
sufficiently large sample counts and on the occupancy event alone,
every allowed individual posterior has excess risk at least $1/100$
for each of the following conventions: posterior-average raw risk,
ordinary raw posterior mean, clipping that mean, posterior-average
clipped risk, and the mean of individually clipped posterior outputs.
The sample probability is at least $1-2e^{-n^{1/3}/2}$.
\end{corollary}
\begin{proof}
In the low-temperature branch, the expected noisy fitting residual
is at most $q+o(1)$. Jensen gives the same residual bound for the
raw mean and the mean of clipped outputs; clipping a raw mean also
decreases the residual. On distinct observed noisy atoms, the vector
of empirical label means has unweighted squared norm at least $S_1$,
while the unweighted fitting residual is no greater than the weighted
one. Since there are $J<2n$ equiprobable noisy atoms, the population
excess for each of these mean predictions has lower limit at least
\[
 \frac p2(\sqrt{pe^{-p}}-\sqrt q)^2>0.037>1/100.
\]
In the high-temperature branch, the expected absolute inner-pair
half-difference is $o(1)$, by the operator-norm estimate. The raw
mean, its clipping, and the mean of clipped outputs therefore each
have signal-component excess at least $(q/2)(1-o(1))^2$, whose limit
is $1/20$. These bounds imply the assertion at sufficiently large
$n$. Jensen gives the corresponding raw and clipped posterior-average
risks. The constants and thresholds are uniform over all allowed
widths, scales, and temperatures.
\end{proof}

For the localized-noise variant, the same argument gives the
individual-posterior excess $p^2e^{-p}/256$. In its low-temperature
branch the lower bound is at least
$(p/2)pe^{-p}(1/\sqrt2-1/4)^2$, and in its high-temperature branch it
is at least $pe^{-p}/128$. Both exceed the stated constant for
$0<p\le9/10$. The mixture statement is weaker numerically but covers
arbitrary averaging of the raw posterior means.

\begin{corollary}[Width alone does not control approximation]
Fix $H\ge1$ and $b>1/[2(H+1)]$. The mixture conclusion of
Theorem~\ref{thm:relu-mixture-failure} holds with a teacher of width
four and coefficient cap $B=n^b$, at learner widths
$N\ge O_H(n\log n)$. It is simultaneous over every
$0<\sigma^2\le1/N$ and every likelihood temperature.
\end{corollary}
\begin{proof}
Use one neuron per group in the tent construction, so
$\delta=2B^{-L}$. The memorizer and sample event are unchanged.
The high-temperature half-difference is now
$O_H(n^{1/2-Lb}\operatorname{polylog}n)=o(1)$.
All other arguments apply without change.
\end{proof}

The coefficients are a necessary part of a quantitative approximation
class: a distribution-uniform width-only oracle guarantee fails even
with constant comparator width. The input dimension here is
$\Theta(\log n)$ and the law varies with sample count. This does not
refute consistency for every fixed law.

\begin{corollary}[No universal joint sample polynomial]
For the unrestricted ReLU Gaussian-prior learner family considered
here, there is no distribution-uniform sample bound polynomial jointly
in input dimension, comparator width, and comparator depth, independent
of arbitrarily large learner width, even with unit-bounded comparator
coefficients and arbitrary data-dependent mixtures of the allowed
prior scales, likelihood temperatures, and widths.
\end{corollary}
\begin{proof}
Suppose a proposed polynomial has width degree at most $D_0$.
Choose fixed hidden depth $H$ and a real $r$ with
$D_0<r<2H$. Take $a=1/r$ and $b=0$ in
Theorem~\ref{thm:relu-mixture-failure}. Its sample count is of order
$m^r$, which eventually exceeds that polynomial because depth is
fixed and dimension grows only logarithmically. Nevertheless the
excess risk stays above a positive constant, contradicting the
claimed guarantee for a smaller fixed accuracy and failure
probability. The same argument applies to strong-prior minimizers.
\end{proof}

This conclusion is about noisy real-valued regression. It does not
settle whether an exponential depth dependence is necessary for
Boolean classification, or prove exponential depth dependence at one
fixed comparator width. Restrictions on prior support or on comparator
stability change the question and may permit better rates.

The positive and negative strict exponents match for actual finite
$\gamma\ge2$, and at all $\gamma\ge1$ with one hidden layer. The
negative theorem includes all $\gamma\ge1$ at every depth. For
multiple hidden layers the matching positive boundary throughout
$1\le\gamma<2$ remains open; equality on the boundary is also open.
All statements concern exact inference and do not provide an efficient
sampling or optimization algorithm.
